% ---------------------------------------------------------------------
\chapter{Sicherheit und Ordnung im Repository}
\label{kap:sicherheit}

\begin{lernziele}
\begin{itemize}[nosep]
  \item wissen, was nicht in ein Repository gehört und warum,
  \item eine \cmd{.gitignore} lesen und schreiben,
  \item richtig reagieren, wenn ein Passwort versehentlich auf GitHub gelandet ist.
\end{itemize}
\end{lernziele}

\section{Was nicht ins Repository gehört}

\begin{center}
\renewcommand{\arraystretch}{1.35}
\begin{tabularx}{\linewidth}{@{}>{\raggedright\arraybackslash}p{4.2cm}>{\raggedright\arraybackslash}X>{\raggedright\arraybackslash}p{3.6cm}@{}}
  \toprule
  \textbf{Was} & \textbf{Warum nicht} & \textbf{Stattdessen} \\
  \midrule
  Passwörter, API-Schlüssel, Tokens & werden mit dem Repository geteilt und bleiben in der Geschichte & Datei \cmd{.env}, in \cmd{.gitignore} \\
  erzeugte Dateien (PDF aus LaTeX, \cmd{__pycache__}, \cmd{build/}) & lassen sich jederzeit neu erzeugen und verursachen unnötige Konflikte & in \cmd{.gitignore} \\
  grosse Dateien (Videos, Datensätze) & Git speichert jede Version vollständig; das Repository wird riesig. GitHub lehnt Dateien über 100\,MB ab & Cloud-Speicher, Git LFS \\
  persönliche Einstellungen, Daten von Programmen & betreffen nur ein Gerät & lokal lassen \\
  urheberrechtlich geschütztes Material Dritter & in öffentlichen Repositorys eine Rechtsverletzung & privates Repository oder Verweis \\
  \bottomrule
\end{tabularx}
\end{center}

\begin{praxis}
Die Faustregel in der Industrie lautet: \textbf{Ins Repository gehört die Quelle, nicht das
Ergebnis.} Die \cmd{.tex}-Datei kommt hinein, das daraus erzeugte PDF nicht. Der
Python-Code kommt hinein, die erzeugten \cmd{.pyc}-Dateien nicht. Jeder kann das Ergebnis
aus der Quelle neu herstellen.
\end{praxis}

\section{Die Datei .gitignore}

Die Datei \cmd{.gitignore} im Hauptordner listet Dateien und Ordner, die Git ignorieren soll.
Sie erscheinen dann nicht in \cmd{git status} und werden von \cmd{git add -A} übergangen.

\begin{lstlisting}[style=ausgabe]
# Python
__pycache__/
*.pyc
.venv/

# LaTeX: PDF entsteht beim Kompilieren neu
build/
*.aux
*.log

# Geheimnisse
.env
\end{lstlisting}

\begin{center}
\renewcommand{\arraystretch}{1.25}
\begin{tabularx}{\linewidth}{@{}>{\ttfamily}l>{\raggedright\arraybackslash}X@{}}
  \toprule
  \normalfont\textbf{Muster} & \textbf{Wirkung} \\
  \midrule
  build/       & der Ordner \cmd{build} an jeder Stelle, mit allem Inhalt \\
  *.log        & alle Dateien mit der Endung \cmd{.log}; \cmd{*} steht für beliebige Zeichen \\
  .env         & genau diese Datei \\
  !wichtig.log & Ausnahme: diese Datei trotz \cmd{*.log} nicht ignorieren \\
  \texttt{\#} Text & Kommentar \\
  \bottomrule
\end{tabularx}
\end{center}

Prüfen, warum eine Datei ignoriert wird:

\begin{lstlisting}[style=befehle]
git check-ignore -v build/ads.pdf      # -v (verbose): nennt Datei und Zeile der Regel
\end{lstlisting}

\begin{fehler}
Eine Datei wird nachträglich in \cmd{.gitignore} eingetragen, erscheint aber weiter in
\cmd{git status}. Ursache: \cmd{.gitignore} wirkt nur auf Dateien, die Git noch nicht kennt.
Eine bereits committete Datei muss zuerst aus dem Index genommen werden:
\cmd{git rm --cached <datei>}. Die Datei bleibt auf der Festplatte erhalten.
\end{fehler}

\section{Wenn ein Geheimnis auf GitHub gelandet ist}

\begin{achtung}
Ein Passwort oder Schlüssel, der einmal auf GitHub war, gilt als \textbf{veröffentlicht}.
Automatische Programme durchsuchen öffentliche Repositorys rund um die Uhr nach solchen
Schlüsseln, oft innerhalb von Minuten nach dem Push. Das Löschen der Datei in einem neuen
Commit genügt nicht: Der Schlüssel bleibt in der Geschichte lesbar.
\end{achtung}

Richtiges Vorgehen, in dieser Reihenfolge:

\begin{enumerate}
  \item \textbf{Den Schlüssel sofort ungültig machen} und einen neuen erzeugen, beim
        jeweiligen Dienst. Das ist der einzige Schritt, der den Schaden wirklich begrenzt.
  \item Die Datei aus dem Repository entfernen und in \cmd{.gitignore} eintragen.
  \item Bei Bedarf die Geschichte bereinigen lassen. Das ist aufwendig und gehört in
        erfahrene Hände.
\end{enumerate}

\section{Öffentlich oder privat?}

Ein \textbf{öffentliches} Repository kann jeder Mensch lesen und kopieren. Es eignet sich als
Portfolio für eigene Arbeiten und wird von Arbeitgebern gerne angesehen. Ein
\textbf{privates} Repository sehen nur eingeladene Personen. Dorthin gehören Unterlagen
Dritter, persönliche Notizen und alles, was nicht für die Öffentlichkeit bestimmt ist.

\begin{erfahrung}
Ein öffentliches GitHub-Profil ist heute oft Teil einer Bewerbung. Saubere Commits,
verständliche README-Dateien und eine nachvollziehbare Geschichte sagen über die Arbeitsweise
mehr aus als jedes Zeugnis.
\end{erfahrung}

\begin{zusammenfassung}
\begin{itemize}[nosep]
  \item Quelle ja, Ergebnis nein. Keine Geheimnisse, keine grossen Dateien.
  \item \cmd{.gitignore} steuert, was Git übersieht; bereits bekannte Dateien mit \cmd{git rm --cached} lösen.
  \item Geleakter Schlüssel: zuerst ungültig machen, dann aufräumen.
\end{itemize}
\end{zusammenfassung}
